A morfologia matemática é uma subárea do processamento de imagens e da 
visão computacional que tem como objetivo estudar e desenvolver transformações
aplicadas a imagens, tendo como principal arcabouço teórico a teoria dos 
conjuntos \cite{najman_graph-based_2014}. 
Conforme sugere o próprio termo \textit{morfologia}, essa área busca analisar
e extrair informações relacionadas às características geométricas dos objetos
presentes nas imagens, como forma, tamanho e estrutura.

Nessa busca pelo desenvolvimento de transformações, o uso de estrutura de dados,
para representar e aplicar transformações em imagens, é extremamente necessário.
Exemplos de estruturas são: quad-trees \cite{tanimoto_hierarchical_1975, hunter_operations_1979},
binary partition tree \cite{salembier_binary_2000}, 
watershed tree \cite{najman_geodesic_1996}

A \textit{component tree} \cite{salembier_antiextensive_1998} é uma das 
estruturas utilizadas para representar imagens. A ideia fundamental por trás 
dessa estrutura consiste em representar uma imagem a partir da ordenação de 
seu conjunto de pixels e de suas componentes conexas. Para isso, é necessário 
definir o conceito de componente conexa em uma imagem 
\cite{rosenfeld_connectivity_1970,perret_inf-structuring_2015} e, de forma 
mais geral, estabelecer como uma imagem pode ser representada por meio de um 
grafo \cite{najman_graph-based_2014}.

A representação de uma imagem como um grafo é realizada de forma direta, 
na qual os vértices correspondem aos pixels da imagem, enquanto as arestas 
são definidas de acordo com um determinado modelo de conectividade 
\cite{rosenfeld_connectivity_1970}.

Conforme dito anteriormente, uma das premissas para a contrução da component
tree é a ordenação do conjunto de pixels. Existem duas formas de ordenar 
os pixels: de maneira ascendente ou de maneira descendente. Essas duas formas
dão origem a duas estrutura de dados duais, a max tree cuja ordenação é 
ascendente, ou seja, a relação de ordenação dos pixels se dá pelo operador
 $\le$ entre os níves de cinza. E a min tree cuja ordenação é descendente,
 ou seja, a relação de ordenaçõs dos pixels se dá pelo operador $\ge$ entre os níveis
 de cinza.

 A component tree originalmente foi desenhada para realizar operações conexas
 com o propósito de resolver tarefas de filtragem e segmentação 
 \cite{salembier_flat_1995,salembier_connected_2009}. Para realizar tais tarefas
 a operação mais comum utilizada é o corte, que pode consistir na remoção de alguns
 nós baseado na otimização de algum critério ou na remoção de alguns nós baseado
 nos atributos contidos nesses nós. 

Tal estrutura é amplamente utilizada devido ao seu baixo custo computacional de
construção. Segundo \cite{najman_building_2006}, uma imagem com $n$ pixels 
pode ter sua \textit{component tree} construída em tempo $O(n\log n)$. 
Além dessa característica, a \textit{component tree} permite extrair
informações morfológicas relevantes das imagens. Por esse motivo, essa 
estrutura tem sido empregada em abordagens de aprendizado profundo para a 
extração de características, conforme apresentado em 
\cite{alves_image_2020,farfan_cabrera_segmentation_2021,kerautret_combining_2023}.

Nesse contexto, uma dificuldade encontrada na utilização de \textit{component trees}
é a atualização dessa estrutura quando a imagem sofre alguma alteração. Isto é,
quando a imagem sofre alguma alteração não existe nenhum algoritmo para atualizar
\textit{component tree} de modo a representar a imagem transformada. Dessa forma,
nesse caso de uso se faz necessário reconstruir a árvore novamente para cada 
mudança que ocorre na imagem.

Em \cite{luz_alves_component_2026} os autore propõem a operação de \textit{merge} (mescla)
de nós para resolver parcialmente o problema de atualização. Como apontado também em 
\cite{luz_alves_component_2026} o problema dual, \textit{split}/divisão, também
precisa ser resolvido para poder-se atualizar completamente a árvore. 

Sob esse contexto, esse projeto de pesquisa pretende contribuir com essa lacuna
na literatura. Para simplificar, assim como \cite{luz_alves_component_2026} trata
apenas de \textit{max tree}, esse projeto tratará também apenas desse tipo de estrutura. 

Este documento está organizado da seguinte forma. Na Seção \ref{sec:objetivos}, 
será apresentada a pergunta de pesquisa, bem como as perguntas auxiliares que a 
orientam. Para isso, serão introduzidas algumas das notações, propriedades e 
definições apresentadas em \cite{luz_alves_component_2026}. Na Seção 
\ref{sec:metodologia}, serão descritos os métodos que serão empregados para 
alcançar os objetivos propostos, bem como os procedimentos adotados para 
garantir a reprodutibilidade da pesquisa. Na Seção 
\ref{sec:plano_de_trabalho_cronograma}, será apresentado o planejamento das 
atividades e a priorização das etapas a serem executadas, de modo a assegurar o 
desenvolvimento adequado da pesquisa. Por fim, na Seção 
\ref{sec:mitigacao_de_riscos}, serão discutidos os possíveis riscos 
associados à pesquisa e propostas estratégias de mitigação e alternativas 
para os casos em que o planejamento originalmente estabelecido não 
produza os resultados esperados.
